\section{Conclusion}
\label{sec:conclusion}

This study asks one bounded question: when should conditional evidence about worker performance change the routing of panoramic layout tasks? The proposed answer is procedural. First separate reference-based accuracy, stable peer compatibility, and worker-attributable structural invalidity. Then estimate conditional risk or failure-family performance using an independent calibration stage, retain only components that cross a frozen support and stability boundary, and freeze the resulting policy before its outcomes are visible. At runtime, unsupported components are disabled and out-of-support decisions revert to calibrated global ranking.

\expectedresult{} If the planned evidence is observed, the final conclusion will be that selected conditional worker-performance components are non-redundant, stable across calibration stages, and useful for routing within their support domain. T1 is expected to show that model assistance can reduce active time while preserving delivery quality, with heterogeneous correction behaviour under higher proposal risk. More importantly, V1 is expected to show that calibration-supported conditional routing passes two ordered safety gates before improving delivery-adjusted quality and resource use relative to global ranking. The corresponding effects, intervals, denominators, and costs remain \tbddata{}.

Until those results exist, the defensible contribution is the auditable evidence contract, not the anticipated direction of the findings. The framework does not claim that task-dependent expertise, adaptive assignment, or fallback was invented here. It specifies when a conditional panoramic-annotation profile is allowed to influence deployment and how that decision must be tested. A final manuscript can claim policy benefit only if the prospective hierarchy supports it; otherwise the correct conclusion is global fallback or policy non-distinguishability.
